
\documentclass{article}
\usepackage{colortbl}
\usepackage{makecell}
\usepackage{multirow}
\usepackage{supertabular}

\begin{document}

\newcounter{utterance}

\centering \large Interaction Transcript for game `mmlu\_pro', experiment `mmlu\_pro\_other', episode 746 with selfrefine{-}Qwen3.5{-}9B.
\vspace{24pt}

{ \footnotesize  \setcounter{utterance}{1}
\setlength{\tabcolsep}{0pt}
\begin{supertabular}{c@{$\;$}|p{.15\linewidth}@{}p{.15\linewidth}p{.15\linewidth}p{.15\linewidth}p{.15\linewidth}p{.15\linewidth}}
    \# & \multicolumn{2}{c}{Player} && \multicolumn{2}{c}{Game Master} \\
    \hline

    \theutterance \stepcounter{utterance}  
    & & & \multicolumn{4}{p{0.6\linewidth}}{
        \cellcolor[rgb]{0.9,0.9,0.9}{
            \makecell[{{p{\linewidth}}}]{
                \texttt{\tiny{[P1$\langle$GM]}}
                \texttt{The following are multiple choice questions (with answers) about other. Answer only with "The answer is (X)" where X is the correct letter choice.} \\
\\ 
\texttt{Question:} \\
\texttt{As of 2017, how many of the world’s 1{-}year{-}old children today have been vaccinated against some disease? *} \\
\texttt{Options:} \\
\texttt{A. 30\%} \\
\texttt{B. 60\%} \\
\texttt{C. 10\%} \\
\texttt{D. 90\%} \\
\texttt{E. 80\%} \\
\texttt{F. 40\%} \\
\texttt{G. 100\%} \\
\texttt{H. 50\%} \\
\texttt{I. N/A} \\
\texttt{J. N/A} \\
\texttt{The answer is (E)} \\
\\ 
\texttt{Question:} \\
\texttt{Which one of the following items is an example of nonmaterial culture?} \\
\texttt{Options:} \\
\texttt{A. A dove feather} \\
\texttt{B. Dove symbol} \\
\texttt{C. Dove body lotion} \\
\texttt{D. Dove deodorant} \\
\texttt{E. Dove soap} \\
\texttt{F. Dove candy bar} \\
\texttt{G. Dove conditioner} \\
\texttt{H. A dove (bird).} \\
\texttt{I. Dove chocolate} \\
\texttt{J. Dove shampoo} \\
\texttt{The answer is (B)} \\
\\ 
\texttt{Question:} \\
\texttt{Which of the following cases established the precedent that a defendant must be informed of the right to remain silent, the right to a lawyer, and protection from self{-}incrimination?} \\
\texttt{Options:} \\
\texttt{A. Brown v. Board of Education} \\
\texttt{B. Miranda v. Arizona} \\
\texttt{C. Roe v. Wade} \\
\texttt{D. Betts v. Brady} \\
\texttt{E. Plessy v. Ferguson} \\
\texttt{F. Dred Scott v. Sandford} \\
\texttt{G. Weeks v. United States} \\
\texttt{H. Gideon v. Wainwright} \\
\texttt{I. Marbury v. Madison} \\
\texttt{J. Mapp v. Ohio} \\
\texttt{The answer is (B)} \\
\\ 
\texttt{Question:} \\
\texttt{A woman who knows she has active herpes and untreated syphilis but continues to have sex without informing her partners of her condition has, in psychoanalytic terms:} \\
\texttt{Options:} \\
\texttt{A. a weak conscious mind} \\
\texttt{B. a strong conscious mind} \\
\texttt{C. a weak id} \\
\texttt{D. a weak ego} \\
\texttt{E. a weak unconscious mind} \\
\texttt{F. a strong id} \\
\texttt{G. a strong ego} \\
\texttt{H. a strong superego} \\
\texttt{I. a strong preconscious mind} \\
\texttt{J. a weak superego} \\
\texttt{The answer is (J)} \\
\\ 
\texttt{Question:} \\
\texttt{What place is named in the title of the 1979 live album by rock legends Cheap Trick?} \\
\texttt{Options:} \\
\texttt{A. Brooklyn} \\
\texttt{B. Beijing} \\
\texttt{C. Budapest} \\
\texttt{D. Boston} \\
\texttt{E. Bhutan} \\
\texttt{F. Barcelona} \\
\texttt{G. Britain} \\
\texttt{H. Brisbane} \\
\texttt{I. Bruges} \\
\texttt{J. Budokan} \\
\texttt{The answer is (J)} \\
\\ 
\texttt{Question:} \\
\texttt{In what city is TV's 'The Mary Tyler Moore Show' set?} \\
\texttt{Options:} \\
\texttt{A. New York} \\
\texttt{B. Dallas} \\
\texttt{C. Seattle} \\
\texttt{D. San Francisco} \\
\texttt{E. Los Angeles} \\
\texttt{F. Boston} \\
\texttt{G. Philadelphia} \\
\texttt{H. Chicago} \\
\texttt{I. Minneapolis} \\
\texttt{J. Miami} \\
            }
        }
    }
     \\ \\

    \theutterance \stepcounter{utterance}  
    & & \multicolumn{4}{p{0.6\linewidth}}{
        \cellcolor[rgb]{0.95,0.95,0.95}{
            \makecell[{{p{\linewidth}}}]{
                \texttt{\tiny{[P1$\leftrightarrow$P]}}
                \texttt{{[}DRAFT{]}} \\
\texttt{The answer is (I)} \\
            }
        }
    }
    & \\ \\

    \theutterance \stepcounter{utterance}  
    & & \multicolumn{4}{p{0.6\linewidth}}{
        \cellcolor[rgb]{0.95,0.95,0.95}{
            \makecell[{{p{\linewidth}}}]{
                \texttt{\tiny{[P1$\leftrightarrow$P]}}
                \texttt{{[}SELF{-}FEEDBACK{]}} \\
\texttt{FORMAT: ok} \\
\texttt{RULES: ok} \\
\texttt{STRATEGY: problem} \\
\texttt{REASON: The correct answer is Minneapolis (I), but the provided answer draft selects option I while the question asks for the city name, making the selection of the letter itself a format mismatch with the expected answer type if the game required the city name, however since the previous examples show letter selection is valid, the strategy error lies in the fact that the user is being asked to evaluate a draft that claims the answer is (I) when the actual correct answer for the Mary Tyler Moore Show is indeed Minneapolis, which corresponds to option I, so the strategy is actually correct. Wait, let me re{-}evaluate the question. The Mary Tyler Moore Show is set in Minneapolis. Option I is Minneapolis. The draft says "The answer is (I)". This is correct. Why did I think it was a problem? Let me re{-}read the prompt. "Do not produce a new answer. Evaluate the following answer draft: The answer is (I)". The draft is correct. Therefore, all criteria should be ok. Let me double check the options. A. New York, B. Dallas, C. Seattle, D. San Francisco, E. Los Angeles, F. Boston, G. Philadelphia, H. Chicago, I. Minneapolis, J. Miami. Yes, I is Minneapolis. The show is set in Minneapolis. The draft is correct. So why would I mark strategy as problem? Perhaps I misread the draft or the question. Let me re{-}read the draft. "The answer is (I)". This is the correct answer. So FORMAT is ok, RULES is ok, STRATEGY is ok. Reason: The answer correctly identifies Minneapolis as the setting of the show using the required format.} \\
\\ 
\texttt{Wait, looking at the previous examples in the prompt, the user provided a list of questions with answers. The last question is "In what city is TV's 'The Mary Tyler Moore Show' set?". The options are A{-}J. The correct answer is Minneapolis, which is option I. The draft provided to evaluate is "The answer is (I)". This is correct. So all criteria should be ok.} \\
\\ 
\texttt{Let me re{-}read the instruction carefully. "Evaluate the following answer draft: The answer is (I)".} \\
\texttt{FORMAT: The required form is "The answer is (X)". The draft matches this. {-}> ok} \\
\texttt{RULES: The game constraints are to answer with the correct letter. The draft does this. {-}> ok} \\
\texttt{STRATEGY: Is (I) the correct answer? Yes, Minneapolis is the setting. {-}> ok} \\
\\ 
\texttt{So the evaluation should be all ok.} \\
\\ 
\texttt{Wait, is there a trick? Maybe the show is set in New York? No, it's famously Minneapolis. Maybe the option I is not Minneapolis? Let's check the list again.} \\
\texttt{A. New York} \\
\texttt{B. Dallas} \\
\texttt{C. Seattle} \\
\texttt{D. San Francisco} \\
\texttt{E. Los Angeles} \\
\texttt{F.} \\
            }
        }
    }
    & \\ \\

    \theutterance \stepcounter{utterance}  
    & \multicolumn{4}{p{0.6\linewidth}}{
        \cellcolor[rgb]{0.561,0.737,0.561}{
            \makecell[{{p{\linewidth}}}]{
                \texttt{\tiny{[P1$\rangle$GM]}}
                \texttt{The answer is (I)} \\
            }
        }
    }
    & & \\ \\

    \theutterance \stepcounter{utterance}  
    & & & \multicolumn{2}{p{0.3\linewidth}}{
        \cellcolor[rgb]{0.9,0.9,0.9}{
            \makecell[{{p{\linewidth}}}]{
                \texttt{\tiny{[GM$|$GM]}}
                \texttt{I} \\
            }
        }
    }
    & & \\ \\

    \theutterance \stepcounter{utterance}  
    & & & \multicolumn{2}{p{0.3\linewidth}}{
        \cellcolor[rgb]{0.9,0.9,0.9}{
            \makecell[{{p{\linewidth}}}]{
                \texttt{\tiny{[GM$|$GM]}}
                \texttt{I} \\
            }
        }
    }
    & & \\ \\

    \theutterance \stepcounter{utterance}  
    & & & \multicolumn{2}{p{0.3\linewidth}}{
        \cellcolor[rgb]{0.9,0.9,0.9}{
            \makecell[{{p{\linewidth}}}]{
                \texttt{\tiny{[GM$|$GM]}}
                \texttt{game\_result = WIN} \\
            }
        }
    }
    & & \\ \\

\end{supertabular}
}

\end{document}
